\documentclass[12pt,a4paper]{article}

\usepackage[utf8]{inputenc}
\usepackage[ngerman]{babel}
\usepackage[T1]{fontenc}
\usepackage{amsmath}
\usepackage{amsfonts}
\usepackage{amssymb}
\usepackage{graphicx}
\usepackage[left=2cm,right=2cm,top=2cm,bottom=2cm]{geometry}
\usepackage{multicol}
\usepackage{booktabs}
\usepackage[hidelinks]{hyperref}
\usepackage{tikz}
\usepackage{pgfplots}
\usepackage{blindtext}
\usepackage{array}
\usepackage{multirow}
\usepackage{bigdelim}
\usepackage{colortbl}
\usepackage{fancyhdr} 
\usepackage{tabularx}
\usepackage{xcolor}
\usepackage{color}
\usetikzlibrary{decorations.text}
\usetikzlibrary{tikzmark}
\pagestyle{fancy} 
	\fancyhf{} 
	\fancyhead[L]{\includegraphics[scale=0.05]{Bilder/dhbw.png}} 
	\fancyhead[C]{\slshape Analysis} 
	\fancyhead[R]{\slshape LaTeX Version}

\usepackage{helvet}
\renewcommand{\familydefault}{\sfdefault}

\author{\slshape Robin Rausch, Florian Maslowski}
\title{Analysis}
\date{\slshape \today}
\begin{document}
\maketitle
\tableofcontents
\newpage

\section{Eigenwerttheorie}
	Der Eigenvektor $\overrightarrow{x} $ ist der Vektor einer Matrix $A$, der sich bei der Multiplikation mit der Matrix
	nur um die Länge mit dem ändert:\\
	\Large $A \cdot \overrightarrow{x} = \lambda \cdot \overrightarrow{x} $\\
	\normalsize sdfs


\section{Quadrik}   
    Die Quadrik ist die Lösungsmenge von quadratischen Gleichungen mit mehreren Variablen.
	
\section{Satz von Bolzano-Weierstrass}
	\begin{enumerate}
		\item Jede beschränkte Folge in $\mathbb{R}$ oder $\mathbb{C}$ hat wenigstens einen Häufungspunkt
  		\item Jede beschränkte Folge in $\mathbb{R}$ oder $\mathbb{C}$ hat wenigstens eine konvergente Teilfolge
	\end{enumerate}
	
\section{Couchikriterium}

\section{Konvergenzkriterium}
	$\lim_{n \to \infty} \; a_n = a \Leftrightarrow \forall_{\varepsilon > 0} \exists _{n_\varepsilon \in  \mathbb{N}}$\\
	$\forall_\epsilon > 0, n \in \mathbb{N}\; \exists\; n_\epsilon: n > n_\epsilon \Rightarrow |g-a_n| < \epsilon $
	Für alle positiven Epsilon und natürliche n, gibt es eine Grenze $n_\epsilon $, nach der alle Folgenglieder um weniger als epsilon vom Grenzwert entfernt sind.\\ 

	\subsection{Satz der monotonen Konvergenz}
		Jede noch monoton wachsende/fallende nach oben/unten beschränkte Folge $(a_n)_{n \in \mathbb{N}}$ ist konvergent.\\

	\subsection{Leibniz-Verfahren}

	\subsection{Sandwich-Prinzip}
		Wenn eine Folge zwischen 2 konvergierenden Folgen mit dem selben Grenzwert liegt, konvergiert diese auch gegen den selben Grenzwert.

	\subsection{Wurzel-/Quotientenkriterium}
		Eine Reihe $ \sum_{}^{} a_k  $ ist absolut konvergent wenn:\\
		$1. \lim_{k \to \infty} \sqrt[k]{|a_k|}  = q < 1$\\
		$2. \lim_{k \to \infty} \frac{|a_{k+1}|}{a_k} $
		
	\subsection{Integralkriterium}
		

	\subsection{Majoranten-/Minorantenkriterium}
		Majorantenkriterium: Die Reihe wird durch eine größere ersetzt, deren Konvergenz bekannt ist.\\
		Minorantenkriterium: Die Reihe wird durch eine kleinere ersetzt, deren Divergenz bekannt ist.

	\subsection{Stetigkeit}
		Eine Funktion ist stetig, falls die Funktion keine Sprünge hat. 
		Linke Seite der Funktion ist gleich rechte Seite der Funktion.\\
		$f(x_0) = \lim_{x \to x_{0^-}} f(x) = \lim_{x \to x_{0^+}} f(x) $ \\

\section{Kurvendiskussion}
	\subsection{Newtonverfahren}
		Annähern an Nullstellen durch Rekursion:\\
		$x_{n+1} = x_n + \frac{f(x_n)}{f´(x_n)}$

\section{Taylor Entwicklung}
	Wird genutzt um aus Funktionen Reihen zu machen.\\ 
	Formel: $f(x) = \sum_{i = 0}^{\infty} \frac{f^{(i)}(x_0)}{i!} \cdot (x-x_0)^i$


	\begin{enumerate}
		\item 
	\end{enumerate}

\end{document}